\subsection{Native Mixed Precision}
FSDP offers a versatile native mixed precision mechanism. In terms of parameter management, it adheres to the standard mixed precision technique, which maintains both low and full precision copies of parameters~\cite{micikevicius2017mixedprecision}. Forward and backward computation use the low precision, and the optimizer step uses full precision. FSDP permits user-specified precisions for parameters, gradient reduction, and non-trainable buffers, each independently if desired.

For $\Psi$ number of parameter elements (\code{torch.numel}), $K_{\text{low}}$ bytes per low precision element, and $K_{\text{full}}$ bytes per full precision element, this approach to mixed precision normally increases the memory overhead from $K_{\text{full}}\Psi$ to $(K_{\text{low}} + K_{\text{full}})\Psi$ due to maintaining both precision copies. However, FSDP can sidestep the problem given our design to always keep each local sharded \code{FlatParameter} in GPU memory and only dynamically allocate the unsharded \code{FlatParameter}. For $N$ \code{FlatParameter}s with numels given by $\psi_1, \dots, \psi_N$, the parameter peak memory contribution for FSDP actually \textit{decreases} from $\frac{K_{\text{full}}}{F}\sum_{i=1}^N \psi_i + K_{\text{full}} \max_{i=1}^N \psi_i$ to $\frac{K_{\text{full}}}{F}\sum_{i=1}^N \psi_i + K_{\text{low}} \max_{i=1}^N \psi_i$ bytes. In other words, FSDP directly reduces the second $K_{\text{full}} \max_{i=1}^N \psi_i$ term to $K_{\text{low}} \max_{i=1}^N \psi_i$.

In contrast to \code{torch.amp.autocast} that performs just-in-time casts at the operator level, FSDP's native mixed precision only incurs a full-to-low-precision cast per \code{FlatParameter} in its pre-forward and, if resharding after forward, its pre-backward. Moreover, FSDP's mixed precision permits running all collectives in the low precision, which saves communication volume.

% \todo[inline]{Quantitatively compare autocast vs FSDP MP?}

Users most commonly choose \code{FP16} or \code{BF16} as the low precision and \code{FP32} as the full precision. \code{FP16}'s smaller dynamic range compared that of \code{FP32} exposes \code{FP16} to greater risk of numeric underflow and overflow. The standard solution includes a gradient scaler~\cite{AMPGradScaler} that scales gradients to a safe magnitude. However, since FSDP shards gradients across ranks, a normal local gradient scaler implementation breaks mathematical equivalence, and instead, FSDP provides its own sharded gradient scaler.